% ============================================================
% 分片：s02g  对应原文 PDF 第 27--28 页
% 负责章节：§2 Twisting & mirror symmetry；2.4.2 $B$-branes and coherent sheaves
% 译者：translator-04   审校状态：初稿
% ============================================================
% [跨页接续: 原文第26页/27页。本分片自第27页 "We describe the construction ..." 起；
%  其前一句 "For example, it is known to recover ..." 起于第26页末、止于第27页，
%  按 SPEC §3 按语义完整取于 s02f 末尾。]

% ==== 原文 PDF 第 27 页 ====

我们在第 3.2.3 节描述膜范畴的 Hochschild 上同调的构造，重点讨论它与拓扑下降（topological descent）的关系。

在进入任何具体讨论之前，我们指出：若两个理论 $\mathcal{T},\widetilde{\mathcal{T}}$ 互为镜像，则（只要相应的扭变是良定义的）立即可得相应的边界条件范畴互换。例如，理论 $\mathcal{T}$ 中拓扑 $A$ 膜的范畴 $\mathbf{Bdy}_A$ 应等价于理论 $\widetilde{\mathcal{T}}$ 中拓扑 $B$ 膜的范畴 $\widetilde{\mathbf{Bdy}}_B$；$\mathbf{Bdy}_B$ 与 $\widetilde{\mathbf{Bdy}}_A$ 同理。Kontsevich 的同调镜像对称猜想 [41] 把这一想法实现为数学上精确的陈述，并包含我们在第 3 节讨论的所有必要的同伦论考量。

\subsubsection{$B$ 膜与凝聚层}

既然已知二维 TQFT 中的边界条件具有范畴结构，我们现在要描述 $B$ 型扭变的 $\mathcal{N}=(2,2)$ 理论中边界条件的范畴。我们聚焦于 $\sigma$ 模型的情形，许多细节留待附录 C。在该附录中，我们通过边界上辅助 $\mathcal{N}=2$ 量子力学系统的耦合，详细描述 Landau--Ginzburg 模型中的（$\tfrac12$-BPS）边界条件（聚焦于平坦靶空间 $\mathcal{X}=\mathbb{C}^{N}$ 的情形），该辅助系统也称为 Chan--Paton 因子（Chan--Paton factors）。

我们的首要目标是描述未扭变理论中与 $Q_B$ 扭变相容的超对称边界条件，即描述 $B$ 膜范畴 $\mathbf{Bdy}_B$ 中的对象。假设规定玻色子 $\phi|,\bar\phi|$ 的边界值位于子流形 $\mathcal{S}\hookrightarrow\mathcal{X}$ 中，该子流形局部由方程 $f_i(\phi|,\bar\phi|)=0$ 给出。为在 $Q_B\sim\bar\partial_{\mathcal{X}}$ 的作用下不变，这些条件必须是全纯约束：$\bar\partial_{\mathcal{X}}f_i(\phi|,\bar\phi|)=0$，从而 $\mathcal{S}$ 是 $\mathcal{X}$ 的全纯子流形。

% [Tutor Note: 物理图像：$B$ 型扭变下的 $1/2$-BPS 边界条件 = 全纯子流形 $\mathcal{S}\subset\mathcal{X}$ 加上其上的全纯矢量丛（或其有限维复形）；这就是「$B$ 膜 ≈ 凝聚层」的几何来源。边界费米多重态为这些数据提供 $\mathbb{Z}$-分次的复形结构与微分 $\delta$，不同的复形给出不同的 $B$ 膜。]
其余场的边界条件由如下要求唯一确定：边界条件保持由 $Q_B$ 与 $Q_B^{\dagger}=Q_++Q_-$ 生成的完整一维 $\mathcal{N}=2$ 超对称代数。例如，若 $\phi^{1}|=\bar\phi^{\bar 1}|=0$，则必须有 $\bar\psi^{\bar 1}_{+}|+\bar\psi^{\bar 1}_{-}|=0$（因为 $Q_B\bar\phi^{\bar 1}=\bar\psi^{\bar 1}\sim\bar\psi^{\bar 1}_{+}+\bar\psi^{\bar 1}_{-}$），且 $\psi^{1}_{+}|+\psi^{1}_{-}|=0$（因为 $Q_B^{\dagger}\phi^{1}=2\ii\rho^{1}_{t}\sim\psi^{1}_{+}+\psi^{1}_{-}$）。特别地，费米子标量 $\bar\psi|$ 与（拉回到边界上的）玻色子 1-形式 $\rho|$ 必须垂直于 $\mathcal{S}$，而 $\bar\zeta|$ 必须与 $\mathcal{S}$ 相切。

% [Tutor Note: 易错预警：$Q_B\sim\bar\partial_{\mathcal{X}}$ 要求边界条件在 $Q_B$ 下不变，故约束函数 $f_i$ 必须全纯（$\bar\partial_{\mathcal{X}}f_i=0$）；若错用共轭荷 $Q_B^{\dagger}$，则会得到相反的反全纯条件。费米子条件中 $\bar\psi|$、$\rho|$ 法向而 $\bar\zeta|$ 切向，可逐项由 $Q_B$ 与 $Q_B^{\dagger}=Q_++Q_-$ 的不变性核对，指标 $1$ 与 $\bar 1$ 的升降不可混用。]

如附录 C 所述，这些全纯约束通过耦合边界费米多重态（boundary Fermi multiplet）自然地实现。除了自然地实现上述全纯约束外，这些边界费米多重态还能以有限维全纯矢量丛的复形来丰富该全纯子流形，即带有微分 $\delta$ 的、有限维 $\mathbb{Z}$-分次（由 $R$ 对称性给出）的全纯矢量丛。非常粗略地说，这些数据%
% ==== 原文 PDF 第 28 页 ====
编码了靶空间 $\mathcal{X}$ 上的一个凝聚层（coherent sheaf）\footnote{稍微严格地说：注意 $\mathcal{X}$ 的每个开集 $D$ 上都有一个全纯函数环 $\mathbb{C}[D]$。凝聚层 $\mathcal{E}$ 是对每个开集 $D$ 指定一个有限生成 $\mathbb{C}[D]$-模 $\mathcal{E}(D)$ 的数据，并要求这些模在开集重叠处相容地粘合。给定矢量丛 $E\to\mathcal{X}$，可得到相应的模为 $D$ 上全纯截面空间的凝聚层 $\mathcal{E}(D)=\Gamma_D(E)$，其模结构由全纯截面乘以全纯函数给出。类似地，复子流形 $\mathcal{S}$ 上的全纯函数也通过沿包含 $\mathcal{S}\hookrightarrow\mathcal{X}$ 拉回的函数作乘法而给出这样的模。}。

% [原文疑误: 脚注中 "there is an ring" 应为 "there is a ring"，仅为英文冠词笔误，不影响数学内容。]
% [背景补充: 凝聚层是「局部由有限生成模给出、并可在开集重叠处相容粘合」的层，是矢量丛（局部自由层）的推广；摩天层与理想层都是凝聚层而非矢量丛，因此 $B$ 膜的范畴比「流形 + 丛」的几何更丰富。]
\begin{equation}
B\in\operatorname{Obj}(\mathbf{Bdy}_B)\;\leftrightarrow\;\text{凝聚层 }\mathcal{E}\in\operatorname{Obj}(\operatorname{Coh}(\mathcal{X})).
\tag{2.44}
\end{equation}

此外，BPS Hilbert 空间（在平坦半世界面的一条半直线上）的态被等同于该凝聚层的（导出）整体截面，即层上同调（sheaf cohomology）$H^{(0,\bullet)}(\mathcal{X},\mathcal{E})$：$Q_B$ 局域化到常值映射，故 BPS 位形由其在边界上的取值决定，且 $Q_B$ 在边界涨落上表现为 Dolbeault 微分。关于凝聚层（以及更一般地矩阵分解）在此背景下的起源的更代数化的处理，见附录 C。

% [Tutor Note: 推导补足：$Q_B$ 局域化到常值映射，故 BPS 态由边界取值决定；边界上 $Q_B$ 的作用恰好是 Dolbeault 微分 $\bar\partial$，于是 BPS Hilbert 空间 = 凝聚层的层上同调 $H^{(0,\bullet)}(\mathcal{X},\mathcal{E})$，即（导出的）整体截面。这里「导出」指用内射消解计算，以容纳非自由、非局部自由的情形。]

既然已知 $B$ 膜由凝聚层标记，接下来要处理的问题是：确定由凝聚层 $\mathcal{E},\mathcal{F}$ 标记的两个边界条件 $B_1,B_2$ 之间插值的边界局域算子空间。这样的边界局域算子应给出从交汇处之前的 BPS 态到其后的 BPS 态的映射，即从一个截面空间（更好地说：底层复形之间的映射）到另一个截面空间的映射。此外，它应与「用全纯函数（即 $H^{(0,0)}(\mathcal{X})$）乘截面」这一操作对易，后者被视为与体局域算子的碰撞；参见图 5。因此我们得出结论：由凝聚层 $\mathcal{E}$ 与 $\mathcal{F}$ 标记的两个边界条件 $B_1,B_2$ 之间界面上的局域算子，可等同于凝聚层的态射。事实上，凝聚层的所有态射——特别是那些具有非平凡 $R$ 荷／上同调次数的态射——都能以这种方式表示（见例如 [10, 第 3 章]）。这些态射上有自然的微分，物理上对应于 $Q_B$ 对这些局域算子的作用；我们将物理算子等同于态射的 $Q_B$-上同调类（也称为扩张（extensions）或导出态射（derived morphisms））

% [Tutor Note: 推导补足：$\Ext$ 是导出范畴中的 Hom，可用内射或投射消解计算；物理上取 $Q_B$-上同调后，态射的同伦类给出 $\Ext^{k}$ 群，其中 $k$ 由 $R$ 荷／上同调次数决定。因此「边界局域算子 = 凝聚层的态射」只有在取上同调之后才严格成立，这正是导出范畴而非朴素范畴出现的原因。]
\begin{equation}
\Hom_{\mathbf{Bdy}}(B_1,B_2)\cong\Ext_{\operatorname{Coh}(\mathcal{X})}(\mathcal{E},\mathcal{F}).
\tag{2.45}
\end{equation}

最后指出，边界上局域算子的碰撞与凝聚层态射的通常复合一致。

作为例子，考虑单个自由手征多重态 $\Phi$ 的 $B$ 型扭变，即以 $\mathbb{C}$ 为靶空间的凯勒 $\sigma$ 模型。$\mathbb{C}$ 上的凝聚层可等同于多项式代数 $\mathbb{C}[\phi]$ 的模。Dirichlet 边界条件可通过与 $J$ 项为 $J=\Phi$ 的边界费米多重态耦合来构造，得到 $\mathbb{C}[\phi]$-模 $\mathbb{C}[\phi,\gamma]$ 的凝聚层／复形，其中 $\gamma$ 的 $R$ 荷为 $-1$，$\phi$ 以乘法作用，微分为 $Q_B\phi=0$、$Q_B\gamma=\phi$。从 $\mathbb{C}[\phi,\gamma]$ 到自身且与「乘以 $\phi$」对易的映射由 $\phi$（乘以 $\phi$）、$\gamma$（乘以 $\gamma$）与 $\partial_\gamma$（对 $\gamma$ 求导）生成：$\Hom_{\mathbb{C}[\phi]\text{-}\mathrm{mod}}(\mathbb{C}[\phi,\gamma],\mathbb{C}[\phi,\gamma])=\mathbb{C}[\phi,\gamma,\partial_\gamma]$。$Q_B$ 对这些局域算子的作用为

% [Tutor Note: 推导补足：$Q_B\gamma=\phi$、$Q_B\phi=0$ 使 $\gamma$ 成为「乘以 $\phi$」的同伦收缩，故该复形的上同调集中在 $\phi=0$ 处、给出 $\mathbb{C}$（即原点的摩天层），这正对应该 Dirichlet 边界条件所支撑的 $B$ 膜。另一方面 $Q_B\partial_\gamma=0$ 可由 $Q_B$ 是（分次）导子、$Q_B\gamma=\phi$ 而 $Q_B\phi=0$ 直接得到。]
\begin{equation}
Q_B\phi = 0 \qquad Q_B\gamma = \phi \qquad Q_B\partial_\gamma = 0,
\tag{2.46}
\end{equation}
